% ============================================================================
%  3.6 IMPLEMENTACIÓN DEL FRONTEND
% ============================================================================

\section{Implementaci\'on del frontend}

\subsection{Componente \texttt{EventRegistrationButton}}

Se cre\'o un componente reutilizable que gestiona el estado del bot\'on
de inscripci\'on:

\begin{verbatim}
function EventRegistrationButton({
  eventId, eventName, userId,
  isRegistering, isCancelling,
  onRegister, onCancel,
}: EventRegistrationButtonProps) {
  const { data: registration, isLoading } =
    useIsRegistered(userId, eventId);
  const isRegistered = !!registration;

  return (
    <TouchableOpacity
      style={[
        styles.registerButton,
        isRegistered && styles.registerButtonActive,
      ]}
      onPress={() => {
        if (isRegistered) {
          onCancel(eventId, eventName);
        } else {
          onRegister(eventId);
        }
      }}
      disabled={isLoadingState}
    >
      <Ionicons
        name={isRegistered
          ? 'checkmark-circle'
          : 'add-circle-outline'}
        size={16}
        color={isRegistered ? colors.white : colors.primary}
      />
      <Text style={...}>
        {isRegistered ? 'Inscrito' : 'Inscribirse'}
      </Text>
    </TouchableOpacity>
  );
}
\end{verbatim}

\subsection{Integraci\'on en el detalle del escenario}

En \texttt{src/app/scenario/[id].tsx} se modificaron las tarjetas de
evento para incluir el bot\'on de inscripci\'on y el conteo de
inscritos:

\begin{verbatim}
{scenario.events.map((event) => {
  const counts = registrationCounts?.[event.id];
  return (
    <View key={event.id} style={styles.eventCard}>
      {/* ... info del evento ... */}
      {counts && (
        <Text style={styles.registrationCount}>
          {counts.confirmed_count} inscrito(s)
        </Text>
      )}
      <View style={styles.eventActions}>
        {user && (
          <EventRegistrationButton
            eventId={event.id}
            userId={user.id}
            onRegister={handleRegister}
            onCancel={handleCancelRegistration}
          />
        )}
      </View>
    </View>
  );
})}
\end{verbatim}

\subsection{Secci\'on ``Mis inscripciones'' en el perfil}

En \texttt{src/app/(tabs)/profile.tsx} se agreg\'o el componente
\texttt{MyRegistrationsSection} que consulta y muestra las
inscripciones del usuario:

\begin{verbatim}
function MyRegistrationsSection({ userId }) {
  const { data: registrations, isLoading } =
    useMyRegistrations(userId);

  if (isLoading) return <ActivityIndicator />;
  if (registrations?.length > 0) {
    return (
      <View style={styles.infoCard}>
        {registrations.map((reg) => (
          <View key={reg.id}>
            <Ionicons name="checkmark-circle"
              color={colors.success} />
            <Text>Inscrito el {reg.registered_at}</Text>
          </View>
        ))}
      </View>
    );
  }
  return (
    <View style={styles.emptyRegistrations}>
      <Text>No tienes inscripciones aún</Text>
    </View>
  );
}
\end{verbatim}

\subsection{Estilos personalizados}

Se agregaron los siguientes estilos al archivo de estilos del
detalle del escenario:

\begin{table}[H]
\centering
\begin{tabular}{@{}ll@{}}
\toprule
\textbf{Estilo} & \textbf{Descripci\'on} \\
\midrule
eventActions & Contenedor flex para botones de acci\'on \\
registrationCount & Texto verde para conteo de inscritos \\
registerButton & Bot\'on con borde azul, fondo blanco \\
registerButtonActive & Bot\'on con fondo azul (inscrito) \\
registerButtonDisabled & Opacidad reducida (cargando) \\
registerButtonText & Texto azul del bot\'on \\
registerButtonTextActive & Texto blanco del bot\'on activo \\
\bottomrule
\end{tabular}
\caption{Estilos del m\'odulo de inscripci\'on}
\end{table}
